% Part: sets-functions-relations
% Chapter: induction
% Section: induction

\documentclass[../../../include/open-logic-section]{subfiles}

\begin{document}

\olfileid{sfr}{ind}{int}

\olsection{પરિચય}
\sourcecorrection{OLMD-347}{આ સ્રોત-ફાઇલ અન્ય ત્રણ-સ્તરીય વિભાગોની જેમ જ સ્થિત છે, પરંતુ મૂળમાં સમાવેશ-વર્ગ સુધીનો સાપેક્ષ માર્ગ એક સ્તર ઓછો છે. તેને યોગ્ય ત્રણ સ્તરનો માર્ગ આપ્યો છે.}

\begin{explain}
આગમન ગણિત અને તર્કશાસ્ત્રમાં વારંવાર વપરાતી સાબિતી-રીત છે.
ગણિતમાં પ્રાકૃતિક સંખ્યાઓ પરનું આગમન તમે જાણતા હશો.
પરંતુ ગણિત અને તર્કશાસ્ત્રમાં અનેક પ્રકારનાં વસ્તુ-ગણો છે
જેમના પર આગમન વાપરી શકાય છે.

સામાન્ય રીતે આગમન કોઈ સાર્વત્રિક દાવો સાબિત કરવાની છૂટ આપે છે:
કોઈ ચોક્કસ પ્રકારની દરેક વસ્તુને અમુક ગુણધર્મ છે તે બતાવવાનું.
ખાસ કરીને સર્વપરિમાણક દાખલ કરવાની સામાન્ય સાબિતી-રીત
અતિશય જટિલ હોય ત્યારે આગમન ઉપયોગી છે.
આગમન ખાસ રીતે રચાતી ગણિતીય વસ્તુઓ પર જ કામ કરે છે:
જો ગણનો દરેક ઘટક મૂળભૂત હોય અથવા મૂળભૂત ઘટકોમાંથી
રચાયેલો હોય, તો તેના પર આગમન વાપરી શકીએ.
વધુ ચોક્કસ રીતે:
\end{explain}

\begin{defn}
ગણ~$S$ ફલન~$f$ હેઠળ \emph{સંવૃત} છે જો અને ફક્ત જો
જ્યારે પણ $x \in S$ હોય ત્યારે $f(x) \in S$ પણ હોય.
\end{defn}

\begin{explain}
દાખલા તરીકે, પ્રાકૃતિક સંખ્યાઓનો ગણ~($\mathbb{N}$)
ઉત્તરગામી ફલન~$f(x) = x+1$ હેઠળ સંવૃત છે:
જો $x$ પ્રાકૃતિક સંખ્યા હોય, તો $x+1$ પણ છે.
પરંતુ પ્રાકૃતિક સંખ્યાઓ પૂર્વગામી ફલન~$f(x) = x-1$
હેઠળ સંવૃત નથી, કારણ કે $0$ પ્રાકૃતિક સંખ્યા છે, પણ $-1$ નથી.
\end{explain}

\begin{defn}
ગુણધર્મ~$P$ ફલન~$f(x)$ \emph{હેઠળ જળવાય છે} એમ કહીએ,
જો $f$ના પ્રદેશની કોઈ પણ વસ્તુ~$a$ માટે
$P(a)$માંથી $P(f(a))$ ફલિત થાય:
જો $a$ને ગુણધર્મ~$P$ હોય, તો $f(a)$ને પણ હોય.
\end{defn}

\begin{explain}
સમ સંખ્યા હોવાનો ગુણધર્મ ફલન~$f(x) = x+2$ હેઠળ જળવાય છે,
કારણ કે $x$ સમ હોય તો $x+2$ પણ સમ છે.
તે ગુણધર્મ ઉત્તરગામી ફલન હેઠળ જળવાતો નથી,
કારણ કે $x$ સમ હોય તો $x+1$ વિષમ છે.

પ્રાકૃતિક સંખ્યાઓ પર આગમન કામ કરે છે કારણ કે શૂન્યથી
ઉત્તરગામી ફલન સાંત વાર લાગુ પાડીને દરેક પ્રાકૃતિક સંખ્યા મળે છે.
જે ગણો પર આગમન વાપરી શકીએ તે બધામાં આ લક્ષણ છે.
\end{explain}

\begin{defn}[$\Nat$ પર આગમન]
જો શૂન્યને ગુણધર્મ~$P$ હોય અને $P$ ઉત્તરગામી ફલન હેઠળ
જળવાતો હોય, તો દરેક પ્રાકૃતિક સંખ્યાને ગુણધર્મ~$P$ છે.
\end{defn}

\begin{ex}
એકથી $n$ સુધીની પ્રાકૃતિક સંખ્યાઓનો સરવાળો
$\frac{n(n+1)}{2}$ છે. \\
\sourcecorrection{OLMD-348}{મૂળમાં પ્રથમ પ્રાકૃતિક સંખ્યાઓનો સરવાળો કહેવાયો છે, જ્યારે આ જ વિભાગ પ્રાકૃતિક સંખ્યાઓની શરૂઆત શૂન્યથી માને છે. દર્શાવેલું સૂત્ર એકથી આપેલી સંખ્યા સુધીના સરવાળાનું છે. તેથી દાવાનો એ અર્થ સ્પષ્ટ કર્યો છે; પ્રદર્શિત સરવાળામાં શૂન્ય ઉમેરવાથી મૂલ્ય બદલાતું નથી.}

\textbf{આધાર-કિસ્સો.} પ્રથમ શૂન્ય સંખ્યાઓનો, એટલે કે ખાલી
સૂચિનો, સરવાળો $0=
\frac{0\cdot 1}{2}$ છે.\\

\textbf{આગમનનું પગલું.} ધારો કે $k$ માટે ગુણધર્મ સાચો છે.
તે $k+1$ માટે પણ સાચો છે તે બતાવીએ.
બીજા શબ્દોમાં, ધારો કે $P(k)$ સાચું છે, એટલે કે
\[ 0 + 1 + \ldots + k = \frac{k(k+1)}{2} \]
બતાવીએ કે $P(k+1)$ સાચું છે, એટલે કે
\[ 0 + 1 + \ldots + k + (k+1) = \frac{(k+1)(k+2)}{2} \]
અહીં $P(k)$ની ધારણા વાપરીએ.

\begin{align*}
0 + 1 + \ldots + k &= \frac{k(k+1)}{2} \tag{ધારણા} \\
0 + 1 + \ldots + k + (k+1) &= \frac{k(k+1)}{2} + (k+1)
\tag{બંને બાજુએ $k+1$ ઉમેરો} \\
&= \frac{k(k+1) + 2(k+1)}{2} \\
&= \frac{k^2 + k + 2k +2}{2} \\
&=\frac{(k+1)(k+2)}{2}
\end{align*}
આમ આગમનનું પગલું પૂરું થાય છે અને દાવો સાબિત થયો.
\sourcecorrection{OLMD-349}{મૂળના વિસ્તરણમાં પ્રથમ પદ રેખીય છે, જ્યારે અગાઉના ગુણાકારને વિસ્તૃત કરતાં વર્ગપદ જોઈએ. તે પદ સુધાર્યું છે; હવે દર્શાવેલી છેલ્લી ગુણનખંડિત અભિવ્યક્તિ સાથે સમાનતા સાચી છે.}
\end{ex}

\begin{explain}
!!{formula}s માટે મૂળભૂત ઘટકો પરમાણ્વીય !!{formula}s છે.
સંયોજકો વડે ઓછાં જટિલ !!{formula}sમાંથી વધુ જટિલ
!!{formula}s રચાય છે. દરેક સંયોજકને સૂત્રોનું
એવું ફલન સમજી શકાય જે એક કે વધુ ઇનપુટમાંથી નવું
!!{formula} આપે છે. દાખલા તરીકે, બે !!{formula}s
$!A$ અને $!B$ને સંયોજનમાં જોડતું ફલન
$\mathcal E _\land (!A,
!B) = !A \land !B$ તરીકે લખી શકાય.
આવાં ફલનોને \emph{સૂત્ર-રચના ફલનો} કહીએ.
!!{formula}sનો ગણ આ ફલનો હેઠળ સંવૃત છે:
જ્યારે $!A$ અને $!B$ એ !!{formula}s હોય, ત્યારે
$\lnot !A$, $!A \land !B$ વગેરે પણ !!{formula}s છે.
\sourcecorrection{OLMD-350}{મૂળમાં આ ભાગ અને આગળની આગમન-રૂપરેખા દરમિયાન સૂત્રનું વિશેષચિહ્ન તૂટેલું છે. તેને સંયોજનના બીજા સૂત્રની જેમ યોગ્ય મોટા અક્ષરવાળું સૂત્ર-ચિહ્ન આપ્યું છે.}
\sourcecorrection{OLMD-351}{મૂળ દરેક સંયોજકને બે સૂત્ર જોડતું ફલન કહે છે, પરંતુ નકાર એક સૂત્ર પર કામ કરે છે. તેથી એક કે વધુ ઇનપુટ જણાવ્યું છે. પહેલાંની એક-ઇનપુટ સંવૃતતા અને ગુણધર્મ-સંરક્ષણની વ્યાખ્યાઓ અનેક ઇનપુટ માટે આ રીતે સમજવી: બધા ઇનપુટ ગણમાં હોય તો પરિણામ ગણમાં હોય, અને બધા ઇનપુટને ગુણધર્મ હોય તો પરિણામને પણ હોય. પરિમાણકના કિસ્સાઓમાં આ શરત દરેક બાંધનાર ચલની પસંદગી માટે જોઈએ.}
\end{explain}

\begin{defn}[!!{formula}s પર આગમન]
જો દરેક પરમાણ્વીય !!{formula}ને ગુણધર્મ~$P$ હોય અને
$P$ એ !!{formula}-રચના ફલનો હેઠળ જળવાતો હોય,
તો દરેક !!{formula}ને ગુણધર્મ~$P$ છે.
\end{defn}

\begin{explain}
આ વ્યાખ્યા !!{formula}s પરની આગમન-સાબિતી માટે રીત સૂચવે છે.
દરેક !!{formula}ને ગુણધર્મ~$P$ છે તે બતાવવાનું હોય,
તો આ પગલાં લો:\\

\textbf{આધાર-કિસ્સો.} $!A$ પરમાણ્વીય !!{formula} હોય.
[\ldots] તેથી $!A$ને ગુણધર્મ~$P$ છે.

\textbf{આગમનનું પગલું.} $!A$ અને $!B$ એ !!{formula}s હોય
અને બંનેને ગુણધર્મ~$P$ હોય.

કિસ્સો $\lnot$: [\ldots] તેથી $\lnot !A$ ને ગુણધર્મ $P$ છે.

કિસ્સો $\land$: [\ldots] તેથી $!A \land !B$ ને ગુણધર્મ $P$ છે.

કિસ્સો $\lor$: [\ldots] તેથી $!A \lor !B$ ને ગુણધર્મ $P$ છે.

કિસ્સો $\lif$: [\ldots] તેથી $!A \lif !B$ ને ગુણધર્મ $P$ છે.

કિસ્સો $\liff$: [\ldots] તેથી $!A \liff !B$ ને ગુણધર્મ $P$ છે.

કિસ્સો $\lforall[]$: [\ldots] તેથી $\lforall[x] !A$ ને ગુણધર્મ $P$ છે.

કિસ્સો $\lexists[]$: [\ldots] તેથી $\lexists[x] !A$ ને ગુણધર્મ $P$ છે. \\

તેથી દરેક !!{formula}ને ગુણધર્મ~$P$ છે.
\end{explain}

\end{document}
